\section{Literature Review}









\subsection{Immigration and the Changing Working-Age Population}

Immigration is the second major demographic process shaping Spain's labour market. Although Spain was historically a source of emigration, the country became a major destination for international migration in the early 2000s. During the pre-crisis expansion, large immigrant inflows altered the demographic structure of the working population \parencite{amuedo2008complements,arango2013exceptional,gonzalez2011very}. According to the National Institute of Statistics of Spain (Instituto Nacional de Estadística, INE), Spain registered 599,074 immigrations in 2008, when immigration reached a pre-crisis peak and immigrants represented around 13.9\% of the working-age population.

\vspace{0.7em}\par
The 2008 financial crisis changed this trajectory. The burden of job losses was disproportionately borne by foreign workers, immigrant employment declined, and the immigration inflow from abroad fell by more than 50\% between 2008 and 2013 \parencite{rodriguez2016labor}. By 2013, the immigrant share of the labour market had fallen to around 11\%, and Spain experienced net migration outflows. After 2014, however, immigration recovered. Foreign arrivals increased again, and despite the COVID-19 health crisis, inflows continued to rise, reaching 1,258,894 immigrants in 2022. These shifts show that immigration is not only a population process but also a central factor in the changing size and age composition of Spain's potential labour supply.

\subsection{Can Immigration Offset Population Ageing?}

The interaction between ageing and immigration has led to growing interest in replacement migration as a strategy for reducing the labour-market effects of demographic ageing. Several studies suggest that migration can help ease labour-supply downturns, support competitiveness, and foster economic growth \parencite{stepanek2022sectoral,okamoto2021immigration,united2000replacement}. In the Spanish case, labour-force growth has been linked not only to age structure but also to the inflow of non-EU immigrants, suggesting that immigration can influence aggregate labour supply alongside demographic ageing \parencite{casares2018labor}.

\vspace{0.7em}\par
However, the compensating role of immigration is not automatic. By expanding the working population, migration can increase public revenues and help offset demographic pressures, but it can also generate additional demand for public services and affect welfare spending \parencite{fiorio2023migration,naumann2021population}. Immigration may therefore produce both fiscal benefits and costs, depending on migrants' age composition, labour-market attachment, and use of public services \parencite{preston2014effect}. In this sense, the key issue is not whether immigration increases population size in general, but whether it increases effective labour supply through age composition, participation, and hours worked.

